\bookchapter{第四章：安裝和管理App——手機的“技能包”}

\chapterintro{App讓手機無所不能，但也要管好它們。就像你會給新員工安排工位、分配許可權一樣，安裝App也需要一套規範流程。}

\booksection{一、官方商店：最安全的“應用超市”}

\begin{itemize}
  \item \textbf{安卓}：Google Play（國外）或各品牌自帶商店（如華為應用市場、小米應用商店）。
  \item \textbf{iOS}：只有 App Store。
  \item \textbf{鴻蒙}：華為應用市場。
\end{itemize}

核心安全意義：經過稽核、更新有保障、支付安全。行動指南：\textbf{優先只用官方商店下載！}

\booksection{二、純淨/防護模式：開啟手機的“防毒牆”}

主流國產品牌（華為、小米、OPPO、vivo、榮耀）都內建了\textbf{純淨模式}或\textbf{防護模式}，這是系統級別的安全開關。\textbf{強烈建議所有使用者都開啟}。

\booksubsection{它到底是什麼？}

純淨模式（也叫“防護模式”）的\textbf{核心功能是阻斷第三方 App 的安裝}：當它開啟時，系統會\textbf{只允許從官方應用商店安裝 App}，從瀏覽器、二維碼、聊天軟體等渠道下載的安裝包會被直接攔截，從源頭杜絕惡意軟體。

\booksubsection{適用人群}

\begin{itemize}
  \item \textbf{老年使用者}：分不清"官方"和"假冒"廣告，純淨模式是最穩妥的護身符。
  \item \textbf{對應用許可權不敏感的使用者}：不需要安裝小眾 App、測試版、外服遊戲。
  \item \textbf{希望提升裝置安全性的所有使用者}：哪怕是發燒友，開著也比不開強——真有臨時需要，再單獨關一下即可。
\end{itemize}

\booksubsection{開啟路徑（以華為/小米為例）}

\begin{itemize}
  \item \textbf{華為/鴻蒙}：[設定] > [系統和更新] >\textbf{[純淨模式]}，選擇"開啟"。
  \item \textbf{小米}：[設定] > [隱私保護] >\textbf{[保護隱私]}或 [設定] > [應用設定] > [應用安裝]，關閉"安裝未知應用"。
  \item \textbf{OPPO/vivo}：[設定] > [安全] >\textbf{[安裝外部來源應用]}，逐個關閉授權。
\end{itemize}

\booksubsection{應用商店的安全檢測機制}

應用商店不是“擺設”，它背後有完整的安全流水線：

\begin{itemize}
  \item \textbf{病毒掃描}：每個 App 上架前會經過多家防毒引擎聯合掃描，識別已知木馬和病毒特徵。
  \item \textbf{許可權稽核}：商店會審查 App 申請的許可權是否"合理"——一個手電筒 App 要通訊錄許可權？直接駁回或下架。
  \item \textbf{開發者認證}：上架需要企業資質、個人身份證或信用卡認證。出問題能找到人，不是"黑 App"。
  \item \textbf{持續監控}：已上架的 App 也會被持續復掃，新發現的惡意行為會被強制下架。
\end{itemize}

\booksection{三、第三方來源與“電子寄生蟲”：風險極高！}

從網頁或不明連結下載App，可能夾帶病毒、竊取隱私、盜刷支付。這裡我們特別提一個新概念：\textbf{“電子寄生蟲”}。

\textbf{什麼是“電子寄生蟲”？}\par

指的是那些\textbf{被植入木馬的“破解版”軟體}——你以為是“省了錢”，實際是“中了招”。常見形式：

\begin{itemize}
  \item "破解版"、"免費版"、"去廣告版"付費軟體；
  \item "刷贊"、"搶紅包"、"遊戲外掛"等灰色工具；
  \item 群聊裡"小廣告"推薦的所謂"內部版"App。
\end{itemize}

這些程式表面能用，背後可能在\textbf{偷通訊錄、讀簡訊、監聽麥克風、後臺挖礦}。純淨模式是防範它們\textbf{第一道防線}——直接不讓你裝。

\booksubsection{第二道防線建議：}

\begin{itemize}
  \item \textbf{定期全盤病毒掃描}：使用手機自帶安全中心（如華為手機管家、小米安全中心）每月掃一次。
  \item \textbf{謹慎授予敏感許可權}：通訊錄、簡訊、位置、相機、麥克風——這五項"高危許可權"必須嚴格審批。
  \item \textbf{關注耗電與流量異常}：手機突然發燙、電量掉得快、流量半夜跑——八成是後臺藏了"寄生蟲"。
\end{itemize}

如果萬不得已要裝第三方 App：

\begin{enumerate}
  \item 確保來源絕對可信。
  \item 安裝前仔細看許可權要求。
  \item 事後用安全App掃描。
\end{enumerate}

終極建議：在系統設定裡，關掉“安裝未知應用”的許可權。

\booksection{四、裝、更新、解除安裝}

\begin{itemize}
  \item \textbf{安裝}：商店裡點"獲取"或"安裝"，裝完首次開啟會問你要許可權。
  \item \textbf{更新}：\textbf{必須更！}主要是修復安全漏洞。建議開啟自動更新。
  \item \textbf{解除安裝}：長按圖示拖到解除安裝區。想清得更徹底，去 [設定] > [應用管理] 裡找到並解除安裝。
\end{itemize}

\booksection{五、許可權管理：你是App的“老闆”}

App要許可權（比如通訊錄、位置）有時合理（地圖要位置），有時不合理（手電筒要通訊錄）。

\begin{itemize}
  \item \textbf{許可權選項}："始終允許"、"僅在使用時允許"、"禁止"、"每次詢問"。
  \item \textbf{管理原則}：
  \begin{enumerate}
    \item 最小化：非必要不給。
    \item 事後審計：定期去許可權設定裡看看誰在"越權"。
    \item 情境化：像打車軟體，只在用時給位置許可權。
  \end{enumerate}
\end{itemize}

\begin{tcbwarn}{重點盯防的許可權}
\begin{itemize}
  \item 通訊錄/簡訊：可能被用來發垃圾簡訊。
  \item 實時位置：可能被追蹤行蹤。
  \item 相機/麥克風：可能被偷拍偷錄。
  \item 檔案訪問：可能讀取你的隱私檔案。
\end{itemize}

\end{tcbwarn}

\booksection{六、軟體安裝包的構成要素}

你有沒有想過，從應用商店下載的那個安裝包，裡面到底裝了什麼？瞭解安裝包的構成，能幫你更好地理解App是怎麼工作的，也能在遇到問題時知道該從哪裡排查。

\booksubsection{安卓APK檔案結構}

APK（Android Package Kit）本質是一個ZIP壓縮包，解壓後可以看到以下核心目錄和檔案：

\begin{itemize}
  \item \textbf{AndroidManifest.xml}：應用的"身份證"和"說明書"。記錄了應用包名、版本號、所需許可權、支援的系統版本、入口Activity等核心資訊。系統安裝應用時首先讀取這個檔案。
  \item \textbf{classes.dex}：應用的核心程式碼。Java/Kotlin原始碼編譯後變成Dalvik位元組碼，存在這個檔案裡。一個應用可能有多個dex檔案（如classes2.dex、classes3.dex）。
  \item \textbf{resources.arsc}：編譯後的資源索引表。記錄了所有資源的ID和對應關係。
  \item \textbf{res/目錄}：資原始檔夾。包含圖片、佈局檔案、字串、顏色定義等。不同解析度的圖片放在不同的drawable資料夾裡（如drawable-hdpi、drawable-xhdpi）。
  \item \textbf{assets/目錄}：原始資原始檔夾。存放不經過編譯的資源，如網頁檔案、字型、音訊等。App執行時可以透過AssetManager讀取。
  \item \textbf{lib/目錄}：本地庫資料夾。存放編譯好的.so檔案（C/C++程式碼編譯產物），按CPU架構分類（如arm64-v8a、armeabi-v7a）。遊戲和需要高效能運算的App通常有這個目錄。
  \item \textbf{META-INF/目錄}：簽名資訊資料夾。包含MANIFEST.MF、CERT.SF、CERT.RSA三個檔案，用於驗證APK的完整性和簽名。
\end{itemize}

\booksubsection{iOS IPA檔案結構}

IPA（iOS App Store Package）同樣是壓縮包，結構略有不同：

\begin{itemize}
  \item \textbf{Payload/目錄}：核心內容所在。裡面有一個.app資料夾，包含應用的所有內容。
  \item \textbf{Info.plist}：應用的配置檔案，類似安卓的AndroidManifest.xml。記錄了應用ID、版本、許可權需求、支援的裝置等。
  \item \textbf{可執行檔案}：與.app資料夾同名的主程式，是編譯後的Mach-O格式二進位制檔案。
  \item \textbf{Frameworks/目錄}：應用依賴的動態庫和框架。
  \item \textbf{資原始檔}：圖片、音訊、影片等資源，通常存放在.app資料夾內。
  \item \textbf{\_CodeSignature/目錄}：程式碼簽名資訊，用於驗證應用的完整性和來源。
\end{itemize}

\booksubsection{簽名驗證機制}

為什麼從官方商店下載的App更安全？簽名驗證是關鍵。

\begin{itemize}
  \item \textbf{安卓簽名}：開發者用私鑰對APK進行簽名，系統安裝時用公鑰驗證。如果APK被篡改（比如植入病毒），簽名就會失效，系統會拒絕安裝。Google Play還要求App Signing，由Google保管簽名金鑰，進一步保護應用安全。
  \item \textbf{iOS簽名}：蘋果的稽核更嚴格。每個IPA必須有蘋果頒發的證書籤名，否則無法安裝。這就是為什麼iOS只能從App Store或企業證書安裝應用。企業證書籤名的應用有效期有限，且蘋果可以隨時吊銷。
  \item \textbf{簽名的作用}：
  \begin{enumerate}
    \item 驗證來源：確認應用確實來自聲稱的開發者。
    \item 驗證完整性：確保應用在傳輸過程中未被篡改。
    \item 防止冒充：惡意應用無法偽裝成其他應用。
  \end{enumerate}
\end{itemize}

\booksubsection{安裝包大小與最佳化}

你可能會發現，應用商店顯示的App大小和安裝後佔用的空間不一樣。這是因為：

\begin{itemize}
  \item \textbf{下載大小}：壓縮後的APK/IPA大小。
  \item \textbf{安裝大小}：解壓後的大小，通常比下載大小大很多。
  \item \textbf{執行時佔用}：App執行時還會產生快取、資料檔案，佔用更多空間。
\end{itemize}

現代App的最佳化手段：

\begin{itemize}
  \item \textbf{App Bundle（安卓）}：Google Play的動態分發技術。使用者下載時只獲取適合自己裝置的資源（如只下載對應螢幕解析度的圖片），大大減少下載大小。
  \item \textbf{App Thinning（iOS）}：類似的技術，蘋果會根據裝置自動裁剪不需要的資源。
  \item \textbf{按需下載}：遊戲等大型應用可以先下載核心包，後續再下載資源包。
\end{itemize}

\begin{tcbinfo}{小知識：為什麼有的App安裝包特別大？}
遊戲類App通常最大，因為包含大量高畫質貼圖、3D模型、音訊檔案。社交類App也比較大，因為要支援各種功能模組。工具類App通常最小。如果你發現某個工具類App特別大，要警惕是否捆綁了不必要的功能。

\end{tcbinfo}

\begin{tcbpractice}{實操環節：審查3個App的許可權}
選3個常用App（如社交、購物、工具類），進入 [設定] 的應用許可權管理頁，逐項審查，把不必要的敏感許可權關掉或改成“僅在使用時允許”。

\end{tcbpractice}
